\section{Alfred Machine IR}
\label{sec:machine-ir}
\label{sec:discovery-a}

\subsection{Peripheral discovery on an Alfred-native node}

\begin{diagram}
   E2B node boots
        |
        v
   Reads onboard config strap / provisioned identity
   (which peripheral front-end variant, which channels populated
    at manufacturing test -- see Section 5.2 provisioning interface)
        |
        v
   Assembles Capability Manifest (below) from the
   Verified Primitive Library entries compiled into its firmware
        |
        v
   Announces manifest over network (authenticated registration
   handshake with the ACC; no unauthenticated broadcast)
        |
        v
   Alfred Central Compute validates signature, registers device
   instance into this vehicle's Machine IR
\end{diagram}

\subsection{Capability manifest format}

The manifest is the atomic unit both discovery paths (native announce, or Option-B graph compilation) ultimately produce. It is deliberately close to the format sketched in the design brief, formalized here as the actual schema:

\begin{lstlisting}[language=alfredspec]
device:
  id: e2b_node_12.actuator_1
  semantic_type: LinearActuator
  display_name: front_left_door_actuator

capabilities:
  - move_to_position:
      args: {position: percent}
  - move_velocity:
      args: {velocity: percent_per_s}
  - stop: {}

telemetry:
  - position:      {unit: mm, rate_hz: 50}
  - velocity:       {unit: mm_per_s, rate_hz: 50}
  - motor_current:  {unit: A, rate_hz: 200}
  - fault_flags:    {type: bitfield, rate_hz: on_change}

constraints:
  max_current_A: 1.8
  max_velocity_mm_s: 120
  min_position_mm: 0
  max_position_mm: 150
  thermal_derate_above_C: 85

interface:
  node: e2b_node_12
  binding:
    pwm_channel: 2
    hall_input: 3
    current_sense_adc: 0
  network_path: {zone: front_left, phy: 10BASE-T1S, segment: fl_body_spur}

provenance:
  origin: native_manifest        # native_manifest | probe_inferred
  primitive_ref: hkg://primitives/actuator/linear/hbridge_hall_v3
  confidence: 1.0                # always 1.0 for native_manifest
  last_verified: 2026-09-08T00:00:00Z
\end{lstlisting}

The \code{provenance} block is what makes the format shared across both options: a native E2B manifest always carries \code{confidence: 1.0} because the compiler that built the node's firmware also emitted this exact record; a probe-inferred device (Section~\ref{sec:knowledge-graph}) carries a confidence below 1.0 and a reference to the supporting evidence in the Vehicle Knowledge Graph, and Alfred's policy layer (Section~\ref{sec:active-control}) uses that confidence value directly to gate what commands are permitted.

\subsection{Formal Machine IR structure}

The Machine IR for a vehicle is a graph, not a flat device list, because devices depend on network paths which depend on physical nodes which depend on power domains:

\begin{center}
\small
\begin{tabularx}{\linewidth}{@{}l X@{}}
\toprule
\textbf{IR entity} & \textbf{Contents} \\
\midrule
\code{Vehicle} & identity, program/platform reference, architecture type (native / discovered / hybrid) \\
\code{PhysicalNode} & E2B node instance, or (Option B) legacy ECU instance; MCU/HW reference into Hardware Knowledge Graph, firmware version, health status \\
\code{NetworkSegment} & PHY type, bitrate, topology (point-to-point/multidrop/bus), gateway relationships \\
\code{Device} & the capability-manifest record above; belongs to exactly one \code{PhysicalNode}, reachable via one or more \code{NetworkSegment} hops \\
\code{ResourceBinding} & the specific timer/DMA/pin/ADC-channel/CAN-ID/LIN-frame slot a Device consumes on its \code{PhysicalNode} --- this is what the Firmware Compiler's resource planner (Section~\ref{sec:firmware-generation}) both consumes (Option A: must not conflict) and, for Option B legacy adapters, treats as read-only truth it must not violate \\
\code{ControlLoop} & placement decision record (central/edge) plus the deadline/jitter/bandwidth factors from Section~\ref{sec:acc} that produced that decision, so the decision is auditable, not implicit \\
\code{Behavior} & application-layer composition referencing Devices via the Semantic API \\
\bottomrule
\end{tabularx}
\end{center}

\subsection{Data architecture: IR vs. Knowledge Graph vs. capture data}

\begin{diagram}
  +------------------------------------------------------------+
  |                 HARDWARE KNOWLEDGE GRAPH                    |
  |   (cross-vehicle, cross-program -- Alfred's compounding      |
  |    asset, see Section 19)                                    |
  |                                                                |
  |   MCU parts | transceiver init sequences | verified primitives|
  |   | previously-seen CAN-ID semantic patterns per OEM platform |
  +---------------------------+----------------------------------+
                               | referenced by
                               v
  +------------------------------------------------------------+
  |                 MACHINE IR  (per vehicle instance)            |
  |   Vehicle -> PhysicalNode -> NetworkSegment -> Device ->      |
  |   ResourceBinding, ControlLoop, Behavior                      |
  +--------------+--------------------------------+--------------+
                 ^                                  ^
                 | compiled from                    | compiled from
        (Option A)|                          (Option B)|
   +---------------+------+                +------------+-----------+
   | Native Capability     |                | Vehicle Knowledge Graph |
   | Manifests (confidence |                | (Section 8) built from  |
   | = 1.0, emitted by the |                | Capture Data + Protocol |
   | Firmware Compiler)    |                | ID + Correlation Engine |
   +------------------------+                +------------+-----------+
                                                            ^
                                                            | raw input
                                              +-------------+-------------+
                                              | Time-synchronized capture  |
                                              | data (frames, events,      |
                                              | electrical measurements)   |
                                              +----------------------------+
\end{diagram}

This separation matters operationally: capture data is large and vehicle-specific and can be discarded/archived after the Knowledge Graph is built from it; the Vehicle Knowledge Graph is medium-sized and vehicle-specific and is retained as the audit trail behind every Machine IR entry's confidence score; the Machine IR is small, per-vehicle, and is what the Runtime and application layer actually query at run time; the Hardware Knowledge Graph is the only one of the four that grows across the whole fleet/company and is never vehicle-specific.
